\section{Big Data e Mídias Sociais em Finanças}
\label{sec:big_data_midias_sociais}

A expansão das plataformas digitais nas últimas duas décadas transformou profundamente a natureza dos dados disponíveis para a pesquisa em finanças. Se, historicamente, os estudos empíricos dependiam de séries temporais de preços, volumes e indicadores contábeis gerados em ciclos regulares, a consolidação das mídias sociais e dos ambientes de comunicação em tempo real produziu um novo tipo de dado --- massivo, contínuo, heterogêneo e predominantemente textual --- cujo potencial informacional para a compreensão do comportamento dos mercados ainda está sendo explorado \cite{kearney2014,chen2014bigdata}. Esta seção caracteriza esse fenômeno sob o conceito de \textit{Big Data}, examina o papel específico do X/Twitter como fonte de dados financeiros e discute os principais desafios metodológicos que emergem do processamento de grandes volumes de texto financeiro em língua portuguesa.

% ============================================================
%  Seção 2.2.1 — Características de Big Data
%               (Volume, Velocidade, Variedade)
%  Capítulo 2 — Fundamentação Teórica
% ============================================================

\subsection{Características de Big Data (Volume, Velocidade, Variedade)}
\label{subsec:big_data_caracteristicas}

O termo \textit{Big Data} ganhou relevância científica e empresarial a partir do trabalho de \citeauthor{laney2001} (\citeyear{laney2001}), analista do META Group, que identificou três dimensões fundamentais impulsionadas pela explosão do comércio eletrônico no início dos anos 2000: \textbf{volume}, \textbf{velocidade} e \textbf{variedade} --- o chamado modelo dos ``3 Vs''. Embora o conceito tenha evoluído desde então, incorporando dimensões adicionais como veracidade e valor \cite{chen2014bigdata}, os três eixos originais continuam sendo a base conceitual mais adotada na literatura para caracterizar fenômenos de dados em larga escala.

De maneira geral, \textit{Big Data} pode ser definido como a capacidade de coletar, armazenar e analisar conjuntos de dados cujo tamanho, ritmo de geração ou complexidade superam as capacidades dos sistemas gerenciadores de banco de dados relacionais (RDBMS) e das arquiteturas convencionais de \textit{data warehouse}, projetados para dados estruturados em volumes previsíveis e processamento em lote \cite{chen2014bigdata}. Para \citeauthor{chen2014bigdata} (\citeyear{chen2014bigdata}), essa definição é operacionalizada precisamente por meio dos ``3 Vs'', que descrevem, respectivamente, a magnitude dos dados produzidos, a rapidez com que precisam ser processados e a heterogeneidade dos formatos em que se apresentam.

\textbf{Volume} refere-se à quantidade absoluta de dados gerados, armazenados e processados. O crescimento exponencial da produção digital --- impulsionado por  dispositivos conectados, plataformas de mídias sociais, sensores e transações eletrônicas --- fez com que a escala de dados relevantes passasse de gigabytes para terabytes e, atualmente, para petabytes e exabytes \cite{chen2014bigdata}. No contexto financeiro, essa dimensão se manifesta, por exemplo, no volume de ordens de compra e venda negociadas eletronicamente em bolsas de valores, nos registros de operações de alta frequência (\textit{high-frequency trading}) e nas centenas de milhões de publicações diárias em plataformas digitais que tangenciam temas econômicos e financeiros \cite{bollen2011,ranco2015}.

\textbf{Velocidade} diz respeito ao ritmo de geração e à necessidade de processamento dos dados em tempo real ou quase real \cite{laney2001}. Em mercados financeiros, a velocidade é uma dimensão especialmente crítica: informações sobre eventos macroeconômicos, decisões de bancos centrais ou crises corporativas se propagam em milissegundos por redes digitais globais, alterando o comportamento dos investidores e os preços dos ativos antes mesmo que analistas humanos possam interpretar os eventos e antes que sistemas convencionais de processamento em lote (\textit{batch processing}) concluam seus ciclos de atualização \cite{chen2014bigdata}. Plataformas como o X (antigo Twitter) exemplificam essa dimensão ao produzirem fluxos contínuos de publicações que, conforme demonstrado por \citeauthor{bollen2011} (\citeyear{bollen2011}), carregam informações sobre o humor coletivo capazes de antecipar movimentos no índice Dow Jones Industrial Average (DJIA). A análise em tempo real desses fluxos exige infraestrutura computacional capaz de ingerir e processar \textit{streams} de dados de maneira contínua, o que distingue esses sistemas das abordagens convencionais de processamento em lote (\textit{batch processing}).

\textbf{Variedade} refere-se à diversidade de formatos, estruturas e fontes dos dados disponíveis. Enquanto sistemas de informação tradicionais lidam predominantemente com dados estruturados --- organizados em tabelas com campos e tipos definidos --- o ecossistema digital contemporâneo produz volumes massivos de dados semiestruturados (como arquivos JSON e XML) e não estruturados (como textos livres, imagens, vídeos e publicações em redes sociais) \cite{laney2001}. No domínio financeiro, essa variedade é particularmente desafiadora: análises quantitativas baseadas em séries temporais de preços e volumes coexistem com informações qualitativas extraídas de notícias, relatórios corporativos, comunicados de política monetária e conversações em plataformas de mídias sociais \cite{kearney2014}. Integrar essas fontes heterogêneas em modelos preditivos unificados constitui um dos principais desafios metodológicos do campo, como discutido nas seções subsequentes deste trabalho.

Além dos três eixos originais, a literatura recente tem proposto extensões ao modelo. A dimensão de \textbf{veracidade} aborda a incerteza inerente à qualidade e à confiabilidade dos dados, tema especialmente relevante em ambientes de mídias sociais, onde informações falsas, rumores e ruído informacional se misturam a sinais genuínos \cite{chen2014bigdata}. A dimensão de \textbf{valor} ressalta que o simples acúmulo de dados não produz \textit{insights} por si só: é necessário transformá-los em conhecimento acionável por meio de técnicas analíticas adequadas \cite{chen2014bigdata}. No contexto desta dissertação, ambas as dimensões são diretamente relevantes, pois a extração de sinais de sentimento de textos financeiros em língua portuguesa demanda tanto mecanismos de filtragem de ruído quanto abordagens analíticas especializadas capazes de produzir um índice de sentimento coerente e interpretável.

A emergência do \textit{Big Data} redefiniu a fronteira do que é computacionalmente viável na pesquisa em finanças comportamentais. Se, por décadas, a análise empírica em finanças esteve restrita a dados estruturados de alta precisão --- preços, volumes, indicadores contábeis --- a disponibilidade de conjuntos massivos de dados textuais abriu uma nova agenda de pesquisa centrada na extração automatizada de sentimento a partir de fontes digitais \cite{kearney2014,antweiler2004}. É nesse contexto que se insere o presente trabalho: a interseção entre o volume, a velocidade e a variedade característicos das publicações financeiras no X/Twitter e as técnicas de Processamento de Linguagem Natural (PLN) discutidas no Capítulo~\ref{Cap:Fundamentacao_Teorica}, especialmente os modelos baseados em \textit{transformers} adaptados ao domínio financeiro em língua portuguesa.


% ============================================================
%  Subseção 2.2.2 — X/Twitter como Fonte de Dados Financeiros
% ============================================================
\subsection{X/Twitter como Fonte de Dados Financeiros}
\label{subsec:twitter_dados_financeiros}

A utilização de plataformas digitais de comunicação como fonte de dados para análise de mercados financeiros precede o surgimento do X/Twitter. O trabalho pioneiro de \citeauthor{antweiler2004} (\citeyear{antweiler2004}) demonstrou que mensagens publicadas em fóruns de discussão financeira --- especificamente no Yahoo Finance e no Raging Bull --- continham informação relevante sobre o comportamento futuro dos preços de ações. Analisando aproximadamente 1,5 milhão de mensagens relativas a 45 empresas listadas no Dow Jones Industrial Average e no Dow Jones Internet Index entre 1999 e 2001, os autores encontraram evidências de que o volume de mensagens ajudava a prever a volatilidade do mercado, ao passo que o conteúdo dessas mensagens apresentava relação estatisticamente significativa com retornos subsequentes. Embora o efeito sobre retornos fosse economicamente modesto, o estudo estabeleceu o precedente metodológico central para a área: textos gerados por usuários em plataformas digitais carregam conteúdo informacional com implicações para a dinâmica dos preços de ativos.

A ascensão do X/Twitter --- plataforma lançada em 2006 e que rapidamente se tornou um dos principais canais de comunicação em tempo real no mundo --- deslocou o foco da pesquisa dos fóruns especializados para as redes sociais de uso geral. A arquitetura da plataforma, baseada em publicações curtas (\textit{tweets}) de até 280 caracteres, cronológicas e públicas, criou um ambiente propício à disseminação veloz de informações financeiras, opiniões de mercado e reações a eventos econômicos \cite{sprenger2014}. Diferentemente dos fóruns de investidores, o X/Twitter agrega uma audiência heterogênea --- analistas institucionais, investidores individuais, jornalistas e o público geral --- o que amplia tanto o volume de dados disponíveis quanto a diversidade de perspectivas capturadas \cite{mao2011}.

O marco central para a consolidação do X/Twitter como objeto de pesquisa em finanças foi o estudo de \citeauthor{bollen2011} (\citeyear{bollen2011}). Os autores analisaram o conteúdo emocional de cerca de 9,7 milhões de \textit{tweets} publicados entre abril e dezembro de 2008, aplicando duas ferramentas de análise de humor: o OpinionFinder, que mede valência positiva e negativa, e o GPOMS (\textit{Google-Profile of Mood States}), que decompõe o estado de humor coletivo em seis dimensões --- Calma, Alerta, Certeza, Vitalidade, Gentileza e Felicidade. Por meio de análise de causalidade de Granger e de uma rede neural \textit{fuzzy} auto-organizável, os autores demonstraram que a dimensão ``Calma'' do GPOMS apresentava causalidade preditiva sobre os valores de fechamento do DJIA com defasagem de dois a seis dias, alcançando acurácia de 87,6\% na previsão de movimentos diários de alta e baixa do índice. Esse resultado forneceu evidência empírica robusta de que o humor coletivo expresso em mídias sociais contém informação antecipatória sobre o comportamento do mercado, conectando a hipótese de influência do sentimento do investidor \cite{baker2007} ao ambiente digital das redes sociais.

Estudos posteriores ampliaram e refinaram essa linha de investigação. \citeauthor{sprenger2014} (\citeyear{sprenger2014}) examinaram especificamente \textit{tweets} com conteúdo financeiro --- identificados pelo uso de \textit{cashtags} (símbolo \$ seguido do \textit{ticker} da ação) --- e encontraram que tanto o sentimento quanto o volume de publicações relacionadas a uma empresa eram preditivos de retornos e volume de negociação nos dias subsequentes, com efeito mais pronunciado para ações de menor capitalização. \citeauthor{mao2011} (\citeyear{mao2011}), por sua vez, compararam o poder preditivo de \textit{tweets} financeiros com o de buscas no Google Trends, concluindo que a combinação das duas fontes superava qualquer uma isoladamente na previsão de retornos do S\&P~500. Já \citeauthor{ranco2015} (\citeyear{ranco2015}) investigaram a relação entre o volume e o sentimento de \textit{tweets} e os retornos das 30 empresas componentes do DJIA, identificando que o impacto do sentimento no Twitter é especialmente significativo nos dias que antecedem e sucedem eventos corporativos relevantes, como divulgação de resultados e anúncios de fusões.

No contexto brasileiro, a adoção do X/Twitter como fonte de dados financeiros apresenta particularidades relevantes para o presente trabalho. O Brasil figura historicamente entre os países com maior base de usuários da plataforma na América Latina, o que garante um volume substancial de publicações em língua portuguesa com conteúdo relacionado ao mercado financeiro e aos índices da B3. Contudo, a predominância de modelos de análise de sentimento treinados em inglês cria uma lacuna metodológica significativa: léxicos e modelos desenvolvidos para o inglês financeiro --- como o dicionário de \citeauthor{loughran2011} (\citeyear{loughran2011}) --- não capturam adequadamente as nuances semânticas e o vocabulário específico do mercado financeiro brasileiro, demandando o desenvolvimento ou adaptação de ferramentas especializadas para o português \cite{santos2023finbert}. É precisamente essa lacuna que motiva a adoção do modelo FinBERT-PT-BR como ferramenta central de análise de sentimento neste trabalho, conforme detalhado na Seção~\ref{sec:analise_sentimento}.


% ============================================================
%  Subseção 2.2.3 — Desafios de Processamento de Dados Textuais Massivos
% ============================================================

\subsection{Desafios de Processamento de Dados Textuais Massivos}
\label{subsec:desafios_textuais}

A disponibilidade de grandes volumes de dados textuais provenientes de plataformas digitais não implica, por si só, a capacidade de extrair deles informação financeiramente relevante. Entre a coleta bruta de publicações e a geração de sinais de sentimento interpretáveis residem desafios técnicos e linguísticos substanciais, cuja superação constitui um dos problemas centrais da área \cite{kearney2014}. Esses desafios podem ser agrupados em quatro dimensões interdependentes: ruído e qualidade dos dados, especificidade do vocabulário financeiro, escassez de recursos para línguas não inglesas e limitações das abordagens de análise disponíveis.

O primeiro e mais imediato desafio é a razão sinal-ruído característica das mídias sociais. Em plataformas como o X/Twitter, apenas uma fração das publicações que mencionam termos financeiros carrega informação genuinamente relevante para o comportamento dos ativos \cite{bollen2011}. O restante é composto por ruído informacional de diversas naturezas: \textit{spam} automatizado, publicações de robôs (\textit{bots}), comentários irônicos ou sarcásticos cuja polaridade real é oposta à aparente, e menções ambíguas em que o mesmo termo assume sentidos distintos conforme o contexto. A ironia e o sarcasmo são particularmente problemáticos para sistemas de análise de sentimento baseados em léxicos, pois invertem a valência das palavras sem alterá-las formalmente \cite{liu2020survey}. Além disso, a linguagem informal predominante nas redes sociais --- repleta de abreviações, neologismos, emojis e erros ortográficos --- exige etapas de pré-processamento robustas que normalizem o texto antes de qualquer análise semântica \cite{sousa2025}.

O segundo desafio diz respeito à especificidade do vocabulário financeiro. Um dos achados mais influentes nessa área foi demonstrado por \citeauthor{loughran2011} (\citeyear{loughran2011}), ao analisar a aplicação do dicionário Harvard General Inquirer --- léxico de propósito geral amplamente utilizado em ciências sociais --- a relatórios anuais de empresas (formulários 10-K). Os autores identificaram que aproximadamente três quartos das palavras classificadas como negativas por esse dicionário não carregavam conotação negativa no contexto financeiro: termos como \textit{``liability''}, \textit{``tax''} e \textit{``depreciation''} são tecnicamente neutros ou simplesmente descritivos no vocabulário contábil-financeiro, mas aparecem como negativos em léxicos de uso geral. Essa discrepância motivou o desenvolvimento do dicionário Loughran-McDonald, específico para o domínio financeiro em inglês, e evidenciou que a transferência direta de ferramentas de PLN de propósito geral para o domínio financeiro produz resultados sistematicamente distorcidos \cite{loughran2011,kearney2014}.

No contexto brasileiro, essa limitação se desdobra em uma terceira dimensão: a escassez de recursos linguísticos especializados para o português financeiro. Enquanto o inglês dispõe de léxicos validados \cite{loughran2011}, corpora anotados e modelos de linguagem pré-treinados em grandes volumes de texto financeiro, o português --- e em particular o vocabulário do mercado financeiro brasileiro --- permanece subrepresentado nos recursos disponíveis \cite{sousa2025}. Os estudos aplicados ao mercado brasileiro têm documentado sistematicamente essa lacuna: \citeauthor{falcao2020} (\citeyear{falcao2020}), ao analisar notícias do Valor Econômico e da Folha de São Paulo entre 2013 e 2019 e relacioná-las ao comportamento do Ibovespa, e \citeauthor{lucas2025} (\citeyear{lucas2025}), ao aplicar os léxicos VADER e HIV a notícias do mercado brasileiro coletadas via Google News API, constataram que dicionários desenvolvidos para o inglês capturam apenas parcialmente a variância dos índices brasileiros, com o percentual explicado pelos léxicos genéricos permanecendo relativamente baixo. A necessidade de recursos especializados para o português financeiro motivou iniciativas como a
expansão automática de léxicos por técnicas de Word2Vec e Informação Mútua Pontual (PMI) \cite{sousa2025} e, mais recentemente, o desenvolvimento do modelo FinBERT-PT-BR, treinado sobre 1,4 milhão de textos financeiros em português \cite{santos2023finbert}.

O quarto desafio refere-se às limitações inerentes às abordagens disponíveis para processamento em escala. As duas grandes famílias de métodos --- abordagens baseadas em léxicos e abordagens baseadas em aprendizado de máquina --- apresentam compromissos distintos entre interpretabilidade, custo computacional e desempenho \cite{kearney2014}. Abordagens léxicas são transparentes e dispensam dados rotulados, mas falham diante de ironia, negação e vocabulário não coberto pelo léxico. Abordagens baseadas em aprendizado de máquina --- especialmente modelos profundos baseados em \textit{transformers} --- alcançam desempenho superior, mas exigem grandes volumes de dados anotados, que são escassos no domínio financeiro em português \cite{liu2020survey,santos2023finbert}. Além disso, o processamento de fluxos contínuos de \textit{tweets} em tempo real impõe restrições de latência que podem inviabilizar modelos computacionalmente custosos sem infraestrutura adequada \cite{silva2025}. A superação desse conjunto de desafios demanda, portanto, não apenas avanços algorítmicos, mas também a construção cumulativa de recursos linguísticos e computacionais específicos para o domínio financeiro em português, agenda à qual o presente trabalho busca contribuir por meio da adoção e avaliação do FinBERT-PT-BR aplicado às publicações do X/Twitter.